\section{Tokenized-to-Latent Transformations}
\label{sec:symbolic_to_latent}

The transformations covered in Section~\ref{sec:symbolic_internal} kept experience within the tokenized layer, where every artifact re-enters the model through tokenization and occupies context budget. This section crosses the carrier-layer boundary: the target is no longer a discrete token sequence but a continuous representation that participates in the model's hidden-state computation, bypassing tokenization. The underlying reuse primitive shifts from retrieve-and-prepend to attend---the experience occupies a slot in the computation that produces the next action rather than in the text that prompts it.

\subsection{Tokenized-to-Latent Compression}
\label{sec:p3_overview}

P3 takes tokenized experience---raw interaction trajectories or any tokenized experiential record derived from them---and re-encodes it into a continuous latent representation carried by the model's hidden-state space. The operation is compression rather than abstraction or formalization: information is preserved in the geometry of a vector space rather than in the order of tokens. The latent representation is consumed through the attend primitive, removing the per-call prefill cost that long-context tokenized experience incurs under retrieve-and-prepend. Every P3 instance currently in the agent-experience-reuse literature requires a trained composer to produce the latent representation, and all are cross-session implementations that persist across inference calls; session-scoped reuse of inference-time artifacts such as KV-cache compression is structurally permissible under the carrier definition but has not yet been deliberately designed as an experience-transformation primitive in the agent literature. The five works in scope differ in what the composer emits, how it is trained, and what training signal drives it. The remainder of this section first unpacks the four implementation forms of the attend primitive and their associated training signals, then turns to multimodal coverage, and closes with a unified discussion of cost and characteristic weaknesses.

\subsubsection{Mechanism}
\label{sec:p3_variants}

\paragraph{Composer-Generated Latent Tokens.}

In composer-generated latent token variants, a learned module takes the agent's current state as input and emits a fixed-length latent token sequence that is concatenated into the model's input embedding stream and attended as hidden-state slots. MemGen's memory weaver takes the agent's reasoning state as stimulus and constructs a latent token sequence as machine-native memory, while a separate memory trigger monitors the reasoning process and decides when explicit memory invocation is warranted; under unsupervised training conditions the resulting latent memory spontaneously differentiates into functionally distinct faculties---planning memory, procedural memory, and working memory---without explicit supervision for that differentiation. LatentMem targets multi-agent settings with a two-stage design: an experience bank stores raw interaction trajectories in lightweight form, and a memory composer synthesizes compact latent memories conditioned on the retrieved experience and the specific agent's context. The composer is trained through Latent Memory Policy Optimization, which propagates task-level optimization signals through the latent memories back to the composer, encouraging it to produce compact, high-utility representations specific to each agent's role. Both works share the composer-generated mechanism but differ in the training objective: MemGen prioritizes a tight interweaving of memory and reasoning within a single agent, while LatentMem prioritizes role-aware customization across agents. The structural precondition of this variant is that the composer itself must be trained, and the quality of the resulting latent representation is coupled to the volume and diversity of the training interactions.

\paragraph{Retrieved Latent Representations.}

In retrieval-based latent variants, pre-computed latent representations from prior interactions are retrieved by similarity and attended directly, without passing through tokenization. LatentEvolve's daytime scaling instantiates this variant: it retrieves historical latent representations from a growing pool and uses them to better guide current LLM reasoning, treating past latents as reusable computational context that augments the current hidden state. LatentEvolve also contains a nighttime scaling phase that trains a refiner over accumulated latent representations, integrating past latent optimizations to improve future daytime retrieval; that phase constitutes a Latent-to-Parametric step belonging to the composite-pipeline analysis later in the survey and is not unfolded here. The structural vulnerability of this variant is representation drift: the attended-to latent remains fixed while the model parameters and input distribution continue to shift, so the alignment between the stored latent and the current hidden-state manifold degrades over repeated use.

\paragraph{Differentiable Continuous-Space Retrieval.}

In differentiable continuous-space retrieval, the separation between retrieval and generation is dissolved: both operate in a shared continuous space, and a single language-modeling loss jointly optimizes retrieval relevance and generation quality. CLaRa instantiates this variant by constructing key-preserving compressed vectors through a data synthesis procedure based on question answering and paraphrase supervision, and then training the reranker and generator end-to-end via a differentiable top-k estimator that lets gradients flow through both modules. The latent representation is not produced by a standalone composer or retrieved from an external pool but is shaped during end-to-end training alongside the downstream task, making the compression itself part of the optimization objective. At a text compression rate of 16, CLaRa represents the most extreme point on the compression axis among P3 works, but that extremity is purchased at the cost of maximal training coupling: the reranker and generator cannot be updated independently, and any change to the task distribution requires retraining the joint pipeline.

\paragraph{Learnable Parameter Slots.}

In the learnable-parameter-slot variant, a small set of trainable parameters serves as a fixed latent reservoir that is attended at inference time, with the parameter values updated during training by distilling information from a tokenized experience pool. The latent branch of Self-Consolidation instantiates this variant: a contrastive reflection strategy first explicitly summarizes error-prone patterns from past interactions, and a self-consolidation mechanism then distills that non-parametric textual experience into compact learnable parameters, enabling the agent to internalize extensive historical experience directly into its latent space. The same work also contains a narrative-to-narrative branch in which the contrastive reflection output is retained as refined textual experience, discussed under the P1 pathway. The structural advantage of this variant is parameter efficiency---the learnable slots constitute a small fraction of the full model weights---but the fixed slot count introduces a structural tension: incorporating new experience overwrites prior experience encoded in the same slots, and expanding capacity by adding more slots increases the hidden-state computation that must be attended on every forward pass.

Across these four forms, the training signal that drives the composer varies systematically. MemGen's memory weaver is trained through supervised fine-tuning with the agent's task performance providing implicit supervision; LatentMem's composer is trained through LMPO with a task-level reward signal; LatentEvolve's daytime composer is trained jointly with the nighttime refiner, and unpacking that joint signal belongs to the composite-pipeline analysis; CLaRa's reranker and generator are trained jointly through a single language-modeling loss; Self-Consolidation's learnable parameters are updated through a distillation loss from the textual experience pool. Despite this variation in signal type, every P3 work in scope requires an explicit training stage---the training requirement for cross-session latent is structural, not incidental. An additional attend variant, external KV lookup, in which latent key-value pairs are stored and retrieved by query without a trained composer, is structurally covered by the attend primitive but has no representative work in the current agent-experience-reuse literature; its absence is registered as a gap in the discussion. The notion of a discrete \emph{product semantic type}---natural for tokenized carriers that classify outputs as reflections, rules, or workflows---does not transfer to P3, since continuous vector representations carry no discrete semantic labels. MemGen's spontaneously evolved functional clusters---planning, procedural, and working memory---approximate post-hoc semantic types, but they are emergent properties of the latent state space observable only after training, not categories the designer can specify in advance.

\subsubsection{Multimodal Realization}
\label{sec:p3_multimodal}

Modality coverage in the current P3 literature is uniform: all five works in scope are textual. Latent representations are not intrinsically modality-specific---the attend primitive operates on the model's hidden-state space regardless of which modality produced the source experience---but composer training dynamics depend on the encoder side, which is modality-conditioned. The five works in scope all use text encoders; extending the same mechanisms to vision encoders or vision-language-action backbones would introduce training dynamics not observed in the current text-only literature. The absence of multimodal P3 works in the agent-experience-reuse literature is therefore not merely a coverage gap but a mechanistic one: jointly compressing visual perception and textual reasoning into a shared continuous space poses alignment difficulties for which the current literature offers no worked solution.

\subsubsection{Discussion}
\label{sec:p3_discussion}

Compared with the simpler tokenized baseline---prepending the same experience into context---P3 purchases two compounding gains. The latent representation is computed once for a stable body of experience and re-attended without re-tokenization on subsequent calls, so the per-call prefill cost that grows with the length of tokenized experience is amortized. When the composer is trained against a downstream task signal, the latent representation can be optimized through gradients in a continuous space, a form of representation learning that no tokenized carrier permits. The cost side is equally clear: training the composer requires compute investment; the absence of a surface form makes the latent carrier uninspectable and uneditable in the way a stored reflection or workflow is; and cross-session composer training requires a corpus of interactions large enough to align the latent representation with the downstream task. The first of these three, editability, is structural---it follows from the defining property of a Latent carrier, and recovering it requires converting the carrier back to tokenized form, which is P7 territory rather than an engineering fix. The second, compute, is partially engineering-tractable through parameter-efficient composer architectures, distillation from larger composers, and batch optimization, though these reduce training cost without removing the need for training. The third, sample complexity, is engineering-tractable at small scale but structural at large scale: the requirement that a latent representation align with a downstream task implies a sample-complexity floor that cannot be removed by architecture alone.

Two weaknesses recur in cross-session P3 deployments, both tied to the fixed-capacity continuous representation. The first is representation drift: as model parameters and input distributions shift over repeated use, a fixed latent representation gradually loses semantic alignment with the experience it was meant to encode. LatentEvolve's nighttime scaling, which integrates past latent optimizations, can be read as a structural response to drift, but drift itself is not reported as an independently measured failure mode in any of the five works---it is a pathology implied by design choices rather than one characterized by experiment. The second is update interference: when a new experience is composed into a shared latent slot or learnable-parameter reservoir, prior experience in the same slot is partially overwritten. This is the latent analog of the memory bloat that P1 systems exhibit, but the mechanism differs: it is not unbounded growth but overwriting within a fixed-capacity vector. The contrastive reflection that Self-Consolidation uses to control what enters the latent slot is one approach to regulating interference at the input side, but whether it controls interference after distillation is not addressed. The common root of both weaknesses is that a fixed-capacity continuous representation cannot simultaneously remain compact and losslessly accumulate an open-ended stream of experience. The absence of session-scoped P3 implementations in the current literature means that the lowest-cost potential mitigation---training-free latent reuse that accepts the session boundary---has not been explored in the agent context, a gap whose significance is taken up among the open problems.

Latent carriers produced along this pathway are utilized at inference time through Attend-via-Latent, catalogued among the utilization modes of \S\ref{sec:cross_pathway:utilization:modes}.

The transformation covered in this section moves experience reuse from the input space to the hidden-state space, replacing retrieve-and-prepend with attend. Every cross-session instance currently in the literature pays a training cost for that move, and the structurally permitted training-free case---session-scoped latent reuse---remains unpopulated in the agent-experience-reuse literature. Section~\ref{sec:symbolic_to_parametric} continues the sweep along the carrier spectrum into the parametric layer, where experience is internalized into network weights themselves.
